\section{What causes it: four ablations}
\label{sec:ablation}

\needsrun{this entire section. Nothing in it has been run. It is the section that
decides whether the paper is a mechanism paper or a simulation study, and it
should be built before anything else on the to-do list. Full specification in
Appendix~\ref{app:ablation}.}

Section~\ref{sec:phase} establishes that the split happens. It does not establish
what causes it. Four conditions, each swept over the same $(d, f_d)$ grid with
everything else held fixed, isolate the mechanism.

\begin{table}[h]
  \centering
  \small
  \begin{tabular}{@{}llll@{}}
    \toprule
    & \textbf{Condition} & \textbf{Prediction} & \textbf{What it establishes} \\
    \midrule
    A0 & Random-label control     & $\Rmc \approx 0$            & correlation, not noise \\
    A1 & Frozen trust sector      & no transition in $\Rcw$     & trust is the carrier, not the readout \\
    A2 & No anti-learning         & no polarisation at any $d$  & the sign reversal is the engine \\
    A3 & Hebbian / Perceptron     & no label-correlated phase   & specific to reliability-weighted inference \\
    \bottomrule
  \end{tabular}
\end{table}

\paragraph{A0 --- the spurious feature must be \emph{correlated}, not merely
present.} Key the shift of Eq.~\ref{eq:shift} to a per-emitter random bit drawn
independently of the label, at the same magnitude $d$. Every agent still encodes
messages badly, and still does so in a source-dependent way; the only thing
removed is the correlation between that dependence and a label shared across
agents. If $\Rmc$ stays at zero, the effect is not a consequence of adding
uninformative features, and the paper's subject is the correlation rather than
the noise. This is the control the source draft gestures at in prose and never
runs, and it is the cheapest of the four.

\paragraph{A1 --- the trust sector carries the split rather than reporting it.}
Freeze $F_\mu = F_V = 0$, holding every $\mu_{J|I}$ at its initialisation. Note
that $\Rmc$ is then trivially near zero and proves nothing --- the informative
measurement is $\Rcw$. If the opinion sector alone still develops a
label-correlated split, learned reliability is merely where the split becomes
visible; if it does not, reliability estimation is load-bearing. We expect the
latter, and the claim in the title depends on it.

\paragraph{A2 --- the sharpest test.} Replace the prefactor in $F_w$,
\begin{equation}
  \bigl(1 - 2\Phi(\hmu)\bigr) \;\longrightarrow\; \max\!\bigl(0,\; 1 - 2\Phi(\hmu)\bigr),
\end{equation}
so that a distrusted source produces \emph{no} update rather than a reversed one.
Nothing else changes: the same Bayesian machinery, the same trust dynamics, the
same perturbation. This removes the effective antiferromagnetic coupling and
should remove polarisation entirely, including the label-free polarisation at
$d = 0$. If it does, the mechanism is a single sign in a single prefactor, which
is about as sharp as a mechanism claim gets --- and it is directly actionable,
because a system designer who declines to anti-learn from unreliable peers pays a
sample-efficiency cost and buys immunity to this failure.

\paragraph{A3 --- it is not learning-from-peers in general.} Run the same
population, the same perturbation and the same grid under a Hebbian update and
under a Perceptron update. The source draft asserts that neither exhibits the
collective behaviour; asserted, this is worth little, and shown, it is the result
that generalises the paper beyond its own model. It says the phenomenon belongs to
rules that maintain and act on a per-source reliability estimate, which is the
class of rules the introduction is about.

\todo{figure: 4 panels $\times$ the $(d, f_d)$ plane, sharing the colour scale of
Fig.~\ref{fig:phase}, with the unablated result repeated as a fifth reference
panel. One figure, one page of text.}

\paragraph{Reading the four together.} A0 says the perturbation must be
label-correlated. A1 says the trust sector is necessary. A2 says a specific sign
in the trust-to-opinion coupling is necessary. A3 says rules without the trust
sector do not do this at all. Together they narrow the cause from ``a complicated
simulation'' to one identified coupling, and they are falsifiable individually.
